% Back matter: Acknowledgments, Ethical Disclosures, Rights and Permissions, Cite
% This Article, Data Availability, and a Stakeholder Register. The first five
% blocks are final content (keep the wording); the Stakeholder Register is a
% bracketed instruction for the future pass.

\phantomsection
\addcontentsline{toc}{section}{Acknowledgments}
\section*{Acknowledgments}
\label{sec:acknowledgments}

The author would like to acknowledge Anthropic for providing access to Claude
Code for the instruction, code, execution, and document preparation process;
OpenAI for providing access to ChatGPT for deep research summaries; and Google
Gemini for additional search assistance.

\phantomsection
\addcontentsline{toc}{section}{Ethical Disclosures}
\section*{Ethical Disclosures}
\label{sec:ethics}

The author of the article declares no competing interests. This is an independent
draft and is not endorsed, sponsored, or approved by any trial sponsor, CRO, site,
IRB, regulator, or medical society, nor by CFR, ICH, or FDA. All supporting
numbers are simulation results, and the Unitree H2-Surgical 1.0 is a clearly
labelled hypothetical 2030 platform.

\phantomsection
\addcontentsline{toc}{section}{Rights and Permissions}
\section*{Rights and Permissions}
\label{sec:rights}

This article is distributed under the terms of the Creative Commons Attribution
4.0 International License (CC BY 4.0), which permits unrestricted use,
distribution, and reproduction in any medium, provided the original author and
source are properly credited, a link to the Creative Commons license is provided,
and any modifications made are indicated. To view a copy of this license, visit
\url{https://creativecommons.org/licenses/by/4.0/}. This draft was adapted from
original CFR documents in the public domain under 17 U.S.C. \S~105.

\phantomsection
\addcontentsline{toc}{section}{Cite This Article}
\section*{Cite This Article}
\label{sec:cite}

Kawchak K. VVUQ Physical AI Oncology Trial Bill. Zenodo. 2026;
\href{https://doi.org/10.5281/zenodo.xxxxxxxx}{10.5281/zenodo.xxxxxxxx}.

\phantomsection
\addcontentsline{toc}{section}{Data Availability}
\section*{Data Availability}
\label{sec:data}

All scaffolding, instructions, and supporting records for this bill are openly
available in the \texttt{kevinkawchak/cancer-automated} repository under
\texttt{papers/VVUQ-03/draft-paper/}, with the supporting VVUQ-01 and VVUQ-02
records under \texttt{papers/VVUQ-01/} and \texttt{papers/VVUQ-02/}, and are
archived on Zenodo. The national platform, USL, and PSL records are in the sibling
repository \texttt{kevinkawchak/physical-ai-oncology-trials}
(\href{https://doi.org/10.5281/zenodo.19244918}{10.5281/zenodo.19244918}). All
data sets are synthetic or illustrative and reproducible from the deterministic
seed 20260525 described in the supporting records; no patient-identifiable
information is included.

\phantomsection
\addcontentsline{toc}{section}{Stakeholder Register}
\section*{Stakeholder Register}
\label{sec:stakeholders}

[Build a short stakeholder register as one left-aligned ragged-right table
(Stakeholder, Interest, Relevant section) covering the groups named in the policy
memorandum: health systems, clinicians, patient and disability-rights groups, AI
developers, payers, device makers, and regulators. Draw the concerns and the
patient-empowerment framing from the national-platform stakeholder-concerns and
patient-instructions material
(\texttt{physical-ai .../final\_paper/sections/discussion.tex},
\texttt{patient\_instructions.tex}) and tie each stakeholder to the bill section
that addresses its interest. Keep it to one table and two sentences of prose;
cite \texttt{kawchak2026national}, \texttt{usc-ada}, and \texttt{usc-titlevi}.]
